\documentclass[11pt]{article}
\usepackage{amsmath,amssymb,amsthm,hyperref}
\newtheorem{proposition}{Proposition}
\newcommand{\Pp}{\mathcal P}
\newcommand{\B}{\mathcal B}
\title{Rarity-weighted pinning: exact profiles versus coefficient-sensitive splitting}
\author{Campaign 3395: Engineer continuation}
\date{November 4, 2026}
\begin{document}
\maketitle

\paragraph{Status.}
BP2 and the unrestricted fixed-factor lower bound remain unproved.
This note constructs a general, coefficient-preserving pin expansion,
but does not prove its required weighted bound. It identifies an
all-size failure of grouping by \emph{exact} compatibility neighborhoods:
arbitrarily small positive coefficients can make that particular
expansion's rarity-weighted budget exponential.
Splitting the offending addition first repairs the same family, with
one whole-output defect charge and the original sampler.
Thus neither existence of a good weighted expansion nor BP2 is refuted.
The general coefficient-sensitive charge is still missing.
Reviewed V14 is unchanged; no novelty or publication qualification is claimed.

\paragraph{Reused facts and scope.}
We use the reviewed actual compulsory cut, Boolean lift, A4/A7,
J2 and J8--J11, and BP1--BP16, with their original meanings.
The independent Reviewer judgments and the Jukna v1 attribution already
in the ledger are preserved; no external result is newly assumed.
Here $N=n-3$, $L=2^{23}$, $\rho=4/L=2^{-21}$,
$h=\lfloor N/2^{150}\rfloor$, $k=\lfloor h/2\rfloor+1$,
and $\ell=\max(k-2,0)$.
The maximum of a polynomial below is over the \emph{whole} one-hot
assignment domain, as in the global coefficient maxima used in BP9/BP15,
not an assertion that every maximizing assignment belongs to $\Delta$.

\section{An exact expansion, without a quantitative claim}
For fixed-before-sampling nonnegative full-group polynomials $P_i$,
each pinned at a base-source set $I_i$, write
\[
 A_i=\max_\nu P_i(X_\nu),\qquad
 {\cal R}_\rho(P)=\sum_i A_i\rho^{|I_i|},\qquad
 Z=\Pp Q_{\rm heavy}-\sum_iP_i.                         \tag{WP1}
\]
BP8 gives $|\Lambda_\Delta(P_i)|\le A_i\rho^{|I_i|}$.
The desired quantitative assertion is therefore
${\cal R}_\rho(P)-\Lambda_\Delta(Z)\le
C b_Q^\alpha+K M_n^d(s+1)^e e^{-cN}$.
There is no polynomial bound on the number of pieces in this assertion.

\begin{proposition}
Every actual J2 cut admits a finite, sampler-independent nonnegative
full-group expansion
\[
                    \Pp Q_{\rm heavy}=\sum_iP_i,       \tag{WP2}
\]
with fixed base-source pins $|I_i|\ge\ell$.
All formal coefficients, including multiplicities from sharing and
repeated powers, are retained. The construction below groups rows by
their exact compatibility neighborhoods and unions their optional pins.
It asserts neither a small circuit for a component nor a bound on (WP1).
\end{proposition}
\begin{proof}
Fix one actual product $A_jV_j$, retaining its original partition
$S_j=G(A_j)$. Use its exact formal monomials, not parse-tree occurrences.
Let $A^+$ consist of rows with at least $k$ optional nonroot sources
outside $S_j$, and let $A^-=A_j-A^+$. Define $V^+$ on $S_j$.
A heavy pair belongs to exactly one of the two branches
$A^+V_j$ and $A^-V^+$; use the original J2 pair filter in both.

First consider rows in $A^+$ with partners in $V_j$.
Discard internally incompatible monomials for purposes of projection
only. For a row $a$, let ${\cal N}(a)$ be the set of internally
compatible partner monomials whose products with $a$ are compatible.
Partition the rows with nonempty neighborhoods by equality of
${\cal N}(a)$. For a class $E$, let $A_E$ be its row subpolynomial
and $V_E$ its common partner subpolynomial, both with original
coefficients. Project the J2-heavy pairs of $A_EV_E$ to obtain $P_E$.
Every compatible heavy pair occurs once, with its product coefficient.

Let $I_E$ be the union of the optional nonroot base-source sets
of all rows of $E$. Each partner meets every group outside $S_j$.
Because it is compatible with every row of $E$, it uses the same
prescribed nonroot variable at each source in this union.
Conflicting prescriptions would make the neighborhood empty.
Thus every nonzero $P_E$ is pinned at $I_E$.
There are at least $k$ optional source groups in each row, and deleting
the two hubs removes at most two of them, so $|I_E|\ge\ell$.
Filtering out some pairs cannot remove these pins.

Repeat with rows in $V^+$ and partners in $A^-$.
Here the compulsory coverage by $A^-$ is at least the original $S_j$,
which is exactly what the reversed argument needs.
The two disjoint branches and then the original cut indices give (WP2).
Repeated occurrences of a source enter a set union once.
All choices are made from the fixed formal factors, before sampling.

An internally incompatible pair or an incompatible cross pair
contributes nothing to projection. Its Booleanized coefficient remains
in $\B(A_jV_j)-\Pp(A_jV_j)$; it is not silently discarded from the
defect. Summing these coefficients uses A7, or the actual mixed
contexts of A4, with no ordinary-derivative substitution.
\end{proof}

Even allowing a signed residual does not make the charge automatic:
\[
 {\cal R}_\rho(P)-\Lambda_\Delta(Z)
 =-\Lambda_\Delta(\Pp Q_{\rm heavy})
   +\sum_i\bigl(A_i\rho^{|I_i|}+\Lambda_\Delta(P_i)\bigr)
 \ \ge\ -\Lambda_\Delta(\Pp Q_{\rm heavy}).              \tag{WP3}
\]
The last summands are nonnegative by BP8. A useful proof must
bound the residual independently, not assume a benefit from rewriting.

\section{Small coefficients can destroy exact-profile compression}
Choose $v$ disjoint pairs of base sources, with controls on a balanced
side $S$ and probes on its complement. At control $i$ use a root variable
$r_i$ and an allowed nonroot variable $x_i$; at probe $i$ use a root
variable $s_i$ and an allowed nonroot variable $z_i$.
All other outgoing variables remain present in the original problem.
Put
\[
 \begin{split}
 A_c&=\prod_{i=1}^v(r_i+x_i z_i),&
 B&=\prod_{i=1}^v(s_i+z_i),&
 Z_0&=\prod_{i=1}^v z_i,\\
 F_c&=\prod_{i=1}^v(r_i+x_i),&
 C_c&=\prod_{i=1}^v\bigl(r_i(s_i+z_i)+x_i z_i\bigr).
 \end{split}                                                   \tag{WP4}
\]
Let $F_*$ be the full product of the other groups, including the hubs.
Then $P=F_*F_cZ_0$ and $C=F_*C_c$ are coefficient subpolynomials of
the original $F$. For $H=F$ or the actual $H=R_n$, and
$0\le\varepsilon\le1/64$, define the formal positive circuit
\[
 Q_H=\bigl(H-(1-\varepsilon)P-\varepsilon C\bigr)
        +F_*A_c\bigl((1-\varepsilon)Z_0+\varepsilon B\bigr).
                                                               \tag{WP5}
\]
The displayed subtraction specifies a nonnegative baseline polynomial;
it is not a subtraction gate.

\begin{proposition}
For (WP5), on the entire original Boolean cube,
\[
 \Pp Q_H=H,\qquad
 H\le Q_H\le(1+\varepsilon)H,\qquad
 b_{Q_H}=\varepsilon(1-2^{-v}).                         \tag{WP6}
\]
The easy-target circuit has $O(M_n+v)$ arithmetic gates.
No polynomial-size baseline is asserted for $H=R_n$.
There is an original gate order with exactly the heavy polynomial
\[
 W=F_*A_{>h}\bigl((1-\varepsilon)Z_0+\varepsilon B\bigr),
 \quad
 A_{>h}=\sum_{|I|>h}
       \prod_{i\notin I}r_i\prod_{i\in I}x_i z_i.        \tag{WP7}
\]
\end{proposition}
\begin{proof}
The baseline is nonnegative because $H\ge F$ coefficientwise and
$(1-\varepsilon)P+\varepsilon C\le F$ coefficientwise.
Booleanization gives $A_cZ_0=F_cZ_0$, while
$\Pp(F_*A_cB)=C$. Therefore
\[
 Q_H(X)=H(X)+\varepsilon F_*(X)(A_c(X)B(X)-C_c(X)).
                                                               \tag{WP8}
\]
This also identifies the entire output defect.

If one of the selected two-label groups is empty, the parenthesized
difference is zero. Otherwise set
\[
 p_i=\frac{r_i+x_i z_i}{r_i+x_i},\qquad
 q_i=\frac{r_i(s_i+z_i)+x_i z_i}
           {(r_i+x_i)(s_i+z_i)} .
 \]
On Boolean inputs $0\le p_i\le1$ and $q_i\ge p_i/2$.
Consequently
\[
 0\le\frac{A_cB-C_c}{F_cB}
       =\prod_i p_i-\prod_i q_i\le1-2^{-v}.
 \]
Adding other available outgoing variables only increases the full $F$;
the numerator still uses just the selected variables.
Equality holds when every $r_i=0$, $x_i=s_i=z_i=1$,
all other labels at the selected sources are off, and each other
group has one available label. This proves the exact value of $b_Q$.
Since $H\ge F$, (WP6) follows, including every group-zero face.

For $H=F$, write the baseline as
$(1-\varepsilon)(F-P)+\varepsilon(F-C)$.
Both complements have positive first-failing-block recurrences.
For $F-P$, a local block excludes the control/probe patterns
$(r_i,z_i)$ and $(x_i,z_i)$; for $F-C$ it excludes those two and
$(r_i,s_i)$. A local complement is a sum of the remaining positive
terms. With extra labels, partition them into the selected labels
and their positive sum, so this still costs $O(M_n+v)$, not a
quadratic enumeration of the outgoing domains.
Shared prefix/suffix products implement the first-failure sum.
For $H=R_n$, a formal monomial-sum baseline suffices for existence
but can be exponential.

Compute the baseline, the complementary rest factors and
$V=F_{*,\overline S}((1-\varepsilon)Z_0+\varepsilon B)$ first,
then $A=F_{*,S}A_c$. Its compulsory set is exactly $S$ and it is
the last balanced gate; its divided context is exactly $V$.
The subsequent nonzero cuts come from the set-multilinear baseline
and have no optional groups. In this first cut there are no optional
groups on the $V$ side, while an $A$ row indexed by $I$ has exactly
$|I|$ optional nonroot probes. This proves (WP7).
Arbitrary positive integer powers of the selected variables in
$A_c,Z_0$ leave these Boolean and projection identities unchanged.
Shared computations of those powers remain actual shared gates.
\end{proof}

\begin{proposition}
Suppose $0<\varepsilon\le1/64$ and $0\le h\le v-2$.
The exact-neighborhood expansion in (WP2), using even its
\emph{maximal} common base-nonroot pin sets, has cost
\[
 {\cal R}_{\rm profile}
       = B_{v,h}(\rho):=\sum_{j=h+1}^v\binom vj\rho^{2j}.          \tag{WP9}
\]
At $\varepsilon=0$ it instead coalesces to one component of cost
$\rho^v$. Thus the positive-$\varepsilon$ cost is independent of
how small the nonzero perturbation coefficient is.
\end{proposition}
\begin{proof}
A row is indexed by the control word $I$. Its partner neighborhood
contains the $Z_0$ branch and precisely the $B$ words whose nonroot
probe set contains $I$. These neighborhoods are distinct for distinct
$I$ as soon as $\varepsilon>0$.
The resulting projected polynomial is
\[
 \left(\prod_{i\notin I}r_i\prod_{i\in I}x_i\right)F_*
 \left((1-\varepsilon)Z_0+
        \varepsilon\prod_{i\in I}z_i
                    \prod_{i\notin I}(s_i+z_i)\right).            \tag{WP10}
\]
Its one-hot maximum is one, attained when all probes use $z_i$.
Its common nonroot base pins are exactly the controls and probes
indexed by $I$, hence $2|I|$ pins, not just the optional $|I|$.
The other controls are roots; every remaining probe can be a root
in its positive perturbation branch. Other base groups have their
full root-containing factors. The hubs are excluded.
Summing (WP10) for $|I|>h$ proves (WP9).
If $\varepsilon=0$, all nonempty row neighborhoods coincide;
their union pins all probes and the rows merge. Since $h\le v-2$,
no control is a common nonroot pin of the merged component.
\end{proof}

Take $v=\lfloor N/(8L)\rfloor$ and admissible $N\ge2^{153}$.
These sources fit on the two balanced sides, and $v\ge N/(16L)$.
Write $u=\rho^2=2^{-42}$ and $p=u/(1+u)\ge2^{-43}$.
The binomial expansion and the elementary Bernoulli lower-tail
bound used in the reviewed sampler proof give
\[
 B_{v,h}(\rho)
   =(1+u)^v\Pr\{\operatorname{Bin}(v,p)>h\}
   \ge\tfrac12 e^{N/2^{70}}.                           \tag{WP11}
\]
Indeed $vp\ge N/2^{70}$, $h\le vp/2$, and the discarded tail is
at most $e^{-vp/8}<1/2$.
Also $\log(1+u)\ge u/2$ and $v u/2\ge N/2^{70}$.
Choosing $\varepsilon_N=1/(64N^2)$ gives polynomial-size $Q_F$
and $b_Q\to0$, while (WP11) grows exponentially.
This refutes the \emph{generic exact-profile budget shortcut}.
It does not refute an actual-target polynomial-in-original-size
bound: the $R_n$ baseline has no small-size guarantee.

\section{The addition-aware expansion repairs this same family}
\begin{proposition}
Split the actual addition in $V$ before taking compatibility profiles.
For either choice of $H$ in (WP5), set
\[
 P_0=(1-\varepsilon)\Pp(F_*A_{>h}Z_0),\qquad
 Z=\varepsilon\Pp(F_*A_{>h}B).                         \tag{WP12}
\]
Then $\Pp W=P_0+Z$, $P_0$ is pinned at the $v$ probe sources,
and, for the unchanged original sampler,
\[
 {\cal R}_\rho(P_0)-\Lambda_\Delta(Z)
 \le(1-\varepsilon)\rho^v+\varepsilon
 \le 2b_Q+\rho^v
 \le 2b_Q+e^{-N/2^{24}}.                               \tag{WP13}
\]
The last inequality uses the all-size parameters of (WP11).
\end{proposition}
\begin{proof}
The identity follows by distributivity at the specified original
addition, with its original coefficients.
The first maximum is at most $1-\varepsilon$, and on every one-hot
assignment $0\le Z\le\varepsilon$: at most one control word survives,
and $A_{>h}\le A_c\le F_c$ on Boolean inputs.
Apply BP8 to the fixed probe set and use
$-\Lambda_\Delta(Z)\le\mathbb E_\Delta Z\le\varepsilon$.
Equation (WP6) gives $\varepsilon\le2b_Q$ for $v\ge1$.
Finally $21\log2\,v\ge(21/2)N/2^{27}\ge N/2^{24}$.
There is no gate-count multiplier on $b_Q$, no relabeling of a
conditional law as $\Delta$, and no independent multiplication of
rarities at overlapping pins.
\end{proof}

The weighted obstruction therefore has an explicit good expansion.
It kills only exact-profile compression as a generic quantitative
argument, not the live rarity-weighted strategy.
Unlike BP15, it concerns a \emph{weighted} cost, and its failure is
caused by coefficients arbitrarily close to zero rather than merely
an exponential component count.

\section{Exact diagnostic and the still-missing general estimate}
The new source \path{rarity_weighted_pinning_check.py} constructs
(WP5) for $v=4,h=1,\varepsilon=1/128$, with shared squared probe
variables and higher powers at controls.
It checks the actual original cut, exact A4/A7 coefficient accounting,
the profile expansion, and the addition-aware identity.
It also checks the $65536$ selected-label Boolean boxes, all $256$
one-hot assignments, the zero-coefficient coalescence, and a second
strictly positive coefficient. A separate new rational shared product
has both heavy directions and internally incompatible monomials in
both factors; its exact projection partition and A4 contexts are checked.
These are new finite diagnostics,
not an admissible-large-graph sampler enumeration.

There is a small exact LP for the displayed two-branch allocation.
Let $a=\rho^v$ and $B=B_{v,h}(\rho)$.
Retain amount $w$ of the common probe-pinned component and amount
$u$ of each perturbation-row component; pay a uniform residual cap $t$.
The constraints are
\[
 0\le\varepsilon-u\le t,\qquad
 0\le1-w-u\le t,\qquad w,u,t\ge0,
 \quad\min(aw+Bu+t).                                   \tag{WP14}
\]
They follow from the actual one-hot values $\varepsilon$ and $1$.
The optimum \emph{within this stated allocation family} is
\[
                (1-\varepsilon)a+\varepsilon\min(B,1). \tag{WP15}
\]
For fixed $u\in[0,\varepsilon]$, minimizing over $w,t$ sets
$w=1-\varepsilon$, $t=\varepsilon-u$, since $0<a<1$.
The remaining objective is
$(1-\varepsilon)a+\varepsilon+u(B-1)$, proving (WP15).
This is not a global optimization over all possible rewrites.
The first seven-variable, four-equation request was rejected before
optimization by the adapter's simplex-basis admission bound.
Its exact payload and structured rejection are retained; they have
no mathematical significance.
The proved elimination above reduces the LP to
$u+z=\varepsilon$, $u,z\ge0$, with objective $(B-1)u$ and the
separate constant offset $(1-\varepsilon)a+\varepsilon$.
The inspected Jacobian rational-LP contract supplies one bounded
typed computation for this reduction at the actual $\rho=2^{-21}$.
Dedicated JSON files preserve both requests and responses, version,
protocol and error status. The new checker independently verifies
primal/dual witnesses for both the original and reduced problems,
retains the objective offset, and binds them to the circuit's pin table.

\paragraph{Mathematical blocker and changed approach.}
In a general actual circuit, no proof identifies addition branches
whose retained parts have small rarity cost and whose remaining
one-hot signed contribution is controlled by one normalized output
defect budget. A4 and A7 count incompatible coefficients exactly;
they do not bound the compatible portion of a discarded branch.
For (WP5), (WP6) supplies that comparison explicitly. It has not
been extended to arbitrary real coefficients, shared contexts,
both optional directions, and unrestricted degree.
Retire only the generic exact-profile-budget route with (WP9)--(WP11).
Continue with coefficient-sensitive splitting at actual additions
and mixed product contexts, not support-only neighborhood equality.
Neither (WP2), the family-specific (WP13), nor the finite LP proves
BP2, supplies a fixed $\kappa_0>1$, or improves V14.
\end{document}
